\documentclass[11pt,a4paper]{article}
\usepackage{rx-programming-style}

\begin{document}

\rxmaketitle{RX C Programming Style}
\rxdocumentmeta{C source code in RX/Raditex applications, libraries, services, embedded software, and utilities}{1.0}{11 September 2026}

\tableofcontents
\clearpage

\section{1. Purpose}\label{1-purpose}

This document defines the C programming style used in RX projects.

It deliberately shares many rules with the RX PHP style:

\begin{itemize}
\tightlist
\item
  clarity before brevity;
\item
  lowercase underscore identifiers;
\item
  descriptive names;
\item
  Ellemtel/Allman-like braces;
\item
  explicit control flow;
\item
  one clear responsibility per module or function;
\item
  no premature abstraction or optimization.
\end{itemize}

C differs from PHP in one especially important area:

\begin{quote}
Resource ownership and lifetime shall be explicit.
\end{quote}

In C, memory, file descriptors, sockets, library objects, database
results, threads, mutexes, and other resources frequently require
explicit release. Correct cleanup is therefore part of the program
design, not an afterthought.

The words \textbf{shall}, \textbf{should}, and \textbf{may} mean
required, strongly recommended, and optional respectively.

\section{2. Language level}\label{2-language-level}

Portable project code should normally target C11 or a later explicitly
selected C standard.

Example:

\begin{lstlisting}[language=make]
CFLAGS += -std=c11
\end{lstlisting}

When GNU extensions are intentionally required:

\begin{lstlisting}[language=make]
CFLAGS += -std=gnu11
\end{lstlisting}

GNU extensions shall not be introduced accidentally.

\section{3. Core principles}\label{3-core-principles}

\subsection{3.1 Clarity before brevity}\label{31-clarity-before-brevity}

Do not compress control flow merely to save lines.

Avoid:

\begin{lstlisting}[language=C]
if (record == NULL) return -1;
\end{lstlisting}

Prefer:

\begin{lstlisting}[language=C]
if (record == NULL)
{
    return -1;
}
\end{lstlisting}

\subsection{3.2 Explicit state and
ownership}\label{32-explicit-state-and-ownership}

A reader should be able to determine:

\begin{itemize}
\tightlist
\item
  who created an object;
\item
  who owns it;
\item
  how long it remains valid;
\item
  who releases it;
\item
  which function releases it;
\item
  whether a pointer is borrowed or transferred.
\end{itemize}

\subsection{3.3 One responsibility per
module}\label{33-one-responsibility-per-module}

A source/header pair should represent one coherent module.

Example:

\begin{lstlisting}
rx-canopen-client.c
rx-canopen-client.h

rx-mqtt-client.c
rx-mqtt-client.h

rx-device-record.c
rx-device-record.h
\end{lstlisting}

\subsection{3.4 Prefer simple control
flow}\label{34-prefer-simple-control-flow}

Straight-line code with explicit tests is normally preferred to clever
expressions. This is especially important around cleanup and error
paths.

\section{4. Files and headers}\label{4-files-and-headers}

Filesystem names shall use lowercase words separated by hyphens.

Correct:

\begin{lstlisting}
rx-canopen-client.c
rx-canopen-client.h
rx-db-support.c
rx-db-support.h
\end{lstlisting}

Source and header basenames should match.

Public declarations belong in header files. Implementation details
should remain in \passthrough{\lstinline!.c!} files whenever possible.

Use uppercase underscore-style include guards:

\begin{lstlisting}[language=C]
#ifndef RX_CANOPEN_CLIENT_H
#define RX_CANOPEN_CLIENT_H

...

#endif
\end{lstlisting}

A \passthrough{\lstinline!.c!} file should normally include:

\begin{enumerate}
\def\labelenumi{\arabic{enumi}.}
\tightlist
\item
  its own header;
\item
  standard C headers;
\item
  system/third-party headers;
\item
  project headers.
\end{enumerate}

Example:

\begin{lstlisting}[language=C]
#include "rx-canopen-client.h"

#include <stdint.h>
#include <stdio.h>
#include <stdlib.h>
#include <string.h>

#include <glib.h>

#include "rx-log.h"
\end{lstlisting}

\section{5. Identifiers}\label{5-identifiers}

Functions and variables shall use lowercase words separated by
underscores.

Correct:

\begin{lstlisting}[language=C]
rx_canopen_client_open()
rx_canopen_client_close()

database_connection
message_length
device_uuid
sample_count
\end{lstlisting}

Avoid camelCase.

Public functions shall normally carry a project/module prefix:

\begin{lstlisting}[language=C]
rx_canopen_open()
rx_canopen_send_sdo()
rx_canopen_close()
\end{lstlisting}

Names shall be descriptive. Single-letter names shall not be used merely
for brevity.

Prefer:

\begin{lstlisting}[language=C]
for (sample_index = 0;
     sample_index < number_of_samples;
     sample_index++)
{
    ...
}
\end{lstlisting}

instead of:

\begin{lstlisting}[language=C]
for (i = 0; i < n; i++)
{
    ...
}
\end{lstlisting}

Established mathematical notation may be used only when the domain
meaning is clear and the short name genuinely improves readability.

Project type names should use lowercase underscore style, commonly with
a \passthrough{\lstinline!\_t!} suffix where appropriate:

\begin{lstlisting}[language=C]
typedef struct rx_canopen_client rx_canopen_client_t;
\end{lstlisting}

Do not hide pointer semantics behind typedefs such as:

\begin{lstlisting}[language=C]
typedef char * string;
\end{lstlisting}

Macros and compile-time constants shall use uppercase underscore style:

\begin{lstlisting}[language=C]
#define RX_CANOPEN_MAX_NODES 127
#define RX_DEFAULT_TIMEOUT_MS 1000
\end{lstlisting}

\section{6. Formatting and braces}\label{6-formatting-and-braces}

RX C uses an Ellemtel/Allman-like style.

Functions:

\begin{lstlisting}[language=C]
int
rx_device_open(const char *device_name)
{
    ...
}
\end{lstlisting}

Control statements:

\begin{lstlisting}[language=C]
if (connection == NULL)
{
    return -1;
}
else
{
    ...
}
\end{lstlisting}

Loops:

\begin{lstlisting}[language=C]
for (sample_index = 0;
     sample_index < number_of_samples;
     sample_index++)
{
    process_sample(samples[sample_index]);
}
\end{lstlisting}

Use four spaces for indentation. Do not use tabs.

Write the pointer marker next to the variable:

\begin{lstlisting}[language=C]
char *buffer;
const char *device_name;
rx_device_t *device;
\end{lstlisting}

Avoid multiple declarations such as:

\begin{lstlisting}[language=C]
char *first, second;
\end{lstlisting}

Prefer:

\begin{lstlisting}[language=C]
char *first;
char *second;
\end{lstlisting}

Use one statement per line.

\section{7. Functions}\label{7-functions}

Function names shall describe the operation.

Avoid:

\begin{lstlisting}[language=C]
do_it()
proc()
handle()
x()
\end{lstlisting}

Prefer:

\begin{lstlisting}[language=C]
rx_canopen_send_sdo()
rx_database_fetch_device()
rx_wave_read_header()
\end{lstlisting}

Pass only information the function needs.

Use \passthrough{\lstinline!const!} for input data that the function
does not modify:

\begin{lstlisting}[language=C]
int
rx_message_send(
    const rx_connection_t *connection,
    const void *data,
    size_t data_size
);
\end{lstlisting}

Output parameters shall have descriptive names and documented ownership.

Example:

\begin{lstlisting}[language=C]
int
rx_device_read_name(
    const rx_device_t *device,
    char **returned_name
);
\end{lstlisting}

The documentation shall state whether
\passthrough{\lstinline!returned\_name!} receives allocated memory and
how it shall be freed.

Return-value conventions shall be consistent and documented.

\section{8. Integer and size types}\label{8-integer-and-size-types}

Use \passthrough{\lstinline!size\_t!} for object sizes, allocation
sizes, and collection lengths where appropriate.

Use fixed-width integer types for protocol and binary-format fields:

\begin{lstlisting}[language=C]
uint8_t
uint16_t
uint32_t
int32_t
\end{lstlisting}

Do not assume that \passthrough{\lstinline!int!} has a protocol-defined
width.

Avoid accidental signed/unsigned mixing and check narrowing conversions.

\section{9. Memory ownership}\label{9-memory-ownership}

This is a central part of the RX C style.

For every dynamically allocated object it shall be clear:

\begin{lstlisting}
who allocates it
who owns it
when ownership changes
who frees it
how it is freed
\end{lstlisting}

Allocation shall be checked:

\begin{lstlisting}[language=C]
buffer = malloc(buffer_size);

if (buffer == NULL)
{
    return -1;
}
\end{lstlisting}

Do not cast \passthrough{\lstinline!malloc()!} in C:

\begin{lstlisting}[language=C]
buffer = malloc(buffer_size);
\end{lstlisting}

Prefer allocation based on the pointed-to object:

\begin{lstlisting}[language=C]
device = calloc(1, sizeof(*device));
\end{lstlisting}

rather than:

\begin{lstlisting}[language=C]
device = calloc(1, sizeof(rx_device_t));
\end{lstlisting}

Use zero initialization when zero is a valid initial state:

\begin{lstlisting}[language=C]
rx_device_t device = {0};
\end{lstlisting}

For large or untrusted counts, check that multiplication used to
calculate an allocation size cannot overflow.

\section{10. realloc}\label{10-realloc}

Never overwrite the only pointer to allocated memory with
\passthrough{\lstinline!realloc()!} before checking the result.

Incorrect:

\begin{lstlisting}[language=C]
buffer = realloc(buffer, new_size);

if (buffer == NULL)
{
    ...
}
\end{lstlisting}

Correct:

\begin{lstlisting}[language=C]
void *new_buffer;

new_buffer = realloc(buffer, new_size);

if (new_buffer == NULL)
{
    ...
}
else
{
    buffer = new_buffer;
}
\end{lstlisting}

\section{11. Matching acquire/release
APIs}\label{11-matching-acquirerelease-apis}

Use the matching release operation for each resource family.

Typical pairs:

\begin{lstlisting}
malloc/calloc/realloc    -> free
g_malloc/g_new/g_strdup  -> g_free
fopen                    -> fclose
open                     -> close
opendir                  -> closedir
PQexec                   -> PQclear
PQconnectdb              -> PQfinish
GObject reference        -> g_object_unref
\end{lstlisting}

Do not mix unrelated allocation/release APIs.

When a pointer variable remains in scope and may otherwise be reused,
setting it to \passthrough{\lstinline!NULL!} after release can improve
safety:

\begin{lstlisting}[language=C]
free(buffer);
buffer = NULL;
\end{lstlisting}

This does not make other dangling aliases safe; ownership must still be
designed correctly.

\section{12. Cleanup and error paths}\label{12-cleanup-and-error-paths}

Cleanup belongs in the design.

A function that can fail after acquiring several resources shall have a
clear cleanup strategy.

A local \passthrough{\lstinline!goto cleanup!} is acceptable in C when
it makes ownership and error handling clearer.

Example:

\begin{lstlisting}[language=C]
int
rx_file_process(const char *filename)
{
    int result;
    FILE *input_file;
    char *buffer;

    result = -1;
    input_file = NULL;
    buffer = NULL;

    input_file = fopen(filename, "rb");

    if (input_file == NULL)
    {
        goto cleanup;
    }

    buffer = malloc(4096);

    if (buffer == NULL)
    {
        goto cleanup;
    }

    /*
     * Processing...
     */

    result = 0;

cleanup:

    free(buffer);

    if (input_file != NULL)
    {
        fclose(input_file);
    }

    return result;
}
\end{lstlisting}

Resources should normally be released in reverse acquisition order.

An error return shall not leak resources acquired earlier in the
function.

\section{13. Pointer rules}\label{13-pointer-rules}

Pointers returned by fallible operations shall be checked before
dereference.

Borrowed pointers shall be documented:

\begin{lstlisting}[language=C]
/*
 * Returns a borrowed pointer owned by `device`.
 * The caller shall not free it.
 */
const char *
rx_device_get_name(const rx_device_t *device);
\end{lstlisting}

Ownership transfer shall be explicit in API documentation.

Use pointer arithmetic only when it clearly improves the implementation
and bounds are obvious.

\section{14. Strings and buffers}\label{14-strings-and-buffers}

A function that writes into caller-provided storage shall receive the
buffer size:

\begin{lstlisting}[language=C]
int
rx_device_format_name(
    const rx_device_t *device,
    char *buffer,
    size_t buffer_size
);
\end{lstlisting}

Avoid unbounded string functions such as:

\begin{lstlisting}[language=C]
strcpy()
strcat()
sprintf()
gets()
\end{lstlisting}

Prefer size-aware functions such as \passthrough{\lstinline!snprintf()!}
or suitable GLib helpers when GLib is already a dependency.

Do not assume a byte buffer is a C string unless null termination is
guaranteed.

Protocol frames and network data are byte arrays unless the protocol
says otherwise.

\section{15. Structures}\label{15-structures}

Initialize structures deliberately.

Designated initializers are encouraged when they improve clarity:

\begin{lstlisting}[language=C]
rx_message_t message = {
    .type = RX_MESSAGE_STATUS,
    .node_id = node_id,
    .value = value
};
\end{lstlisting}

Private structure definitions should remain in the
\passthrough{\lstinline!.c!} file when callers do not need the layout.

Opaque object-style interfaces are encouraged for substantial modules:

\begin{lstlisting}[language=C]
typedef struct rx_canopen_client rx_canopen_client_t;
\end{lstlisting}

\section{16. Object-like C APIs}\label{16-object-like-c-apis}

C is not object-oriented in the PHP sense, but a module may have an
explicit object lifecycle:

\begin{lstlisting}[language=C]
rx_canopen_client_t *
rx_canopen_client_new(void);

int
rx_canopen_client_open(
    rx_canopen_client_t *client,
    const char *interface_name
);

void
rx_canopen_client_close(
    rx_canopen_client_t *client
);

void
rx_canopen_client_free(
    rx_canopen_client_t *client
);
\end{lstlisting}

The lifecycle should be obvious:

\begin{lstlisting}
new
open
use
close
free
\end{lstlisting}

Do not emulate C++ with complicated inheritance-like macro systems.

\section{17. GLib and GObject}\label{17-glib-and-gobject}

When a module uses GLib, follow GLib ownership conventions consistently.

Typical tools include:

\begin{lstlisting}
g_new0()
g_strdup()
g_free()
GPtrArray
GHashTable
GError
g_object_unref()
\end{lstlisting}

Do not cross allocator families without a documented reason.

Automatic cleanup helpers such as \passthrough{\lstinline!g\_autoptr!}
may be used when supported by the project toolchain and when they make
lifetime easier to understand.

\section{18. Other resources}\label{18-other-resources}

Memory is only one kind of resource.

The same ownership discipline applies to:

\begin{itemize}
\tightlist
\item
  \passthrough{\lstinline!FILE *!};
\item
  POSIX file descriptors;
\item
  sockets;
\item
  \passthrough{\lstinline!DIR *!};
\item
  PostgreSQL \passthrough{\lstinline!PGconn *!};
\item
  PostgreSQL \passthrough{\lstinline!PGresult *!};
\item
  GLib/GObject references;
\item
  parser objects;
\item
  mutexes;
\item
  threads;
\item
  dynamically loaded handles.
\end{itemize}

Every successful acquisition path shall have a matching release path.

\section{19. Error handling}\label{19-error-handling}

Do not ignore return values from operations that can fail when success
is required.

For system calls that use \passthrough{\lstinline!errno!}, preserve or
report it before another call can overwrite it.

Logs should include enough context to identify the failed operation and
object, but shall not expose passwords or secret material.

\section{20. Macros}\label{20-macros}

Prefer functions to function-like macros when a function is sufficient.

Macros are appropriate for compile-time constants, conditional
compilation, and carefully designed low-level helpers.

Macros shall not hide significant control flow or resource ownership.

\section{21. Global state}\label{21-global-state}

Avoid writable global variables.

Prefer explicit module state or object-like context structures passed to
functions.

If global state is unavoidable, document its initialization, lifetime,
and thread-safety.

\section{22. Concurrency}\label{22-concurrency}

Shared mutable state shall have an explicit synchronization policy.

Document:

\begin{itemize}
\tightlist
\item
  which mutex protects which state;
\item
  lock ordering;
\item
  ownership of mutexes;
\item
  whether callbacks may run while locks are held.
\end{itemize}

Do not hold locks across slow I/O unless necessary and documented.

\section{23. Networking and binary
protocols}\label{23-networking-and-binary-protocols}

Use fixed-width integer types for protocol fields.

Do not rely on host byte order.

Validate received lengths before reading fields.

Never trust a network-provided count or length before checking it
against the actual received buffer and project limits.

\section{24. Database C code}\label{24-database-c-code}

For libpq:

\begin{itemize}
\tightlist
\item
  check \passthrough{\lstinline!PQstatus()!} after connecting;
\item
  check \passthrough{\lstinline!PQresultStatus()!} after commands;
\item
  call \passthrough{\lstinline!PQclear()!} for every owned
  \passthrough{\lstinline!PGresult *!};
\item
  call \passthrough{\lstinline!PQfinish()!} for every owned
  \passthrough{\lstinline!PGconn *!};
\item
  use parameterized queries for data values;
\item
  keep connection ownership clear.
\end{itemize}

A cleanup section is often appropriate for functions owning several
libpq resources.

\section{25. Compiler warnings}\label{25-compiler-warnings}

Warnings shall be enabled and treated seriously.

A useful baseline is:

\begin{lstlisting}[language=make]
CFLAGS += -Wall -Wextra -Wformat=2
\end{lstlisting}

Projects should add stricter warnings where practical.

Warnings shall not be disabled merely to get a quiet build without
understanding their cause.

A warning-free normal build is the target.

\section{26. Runtime checking}\label{26-runtime-checking}

Memory-sensitive code should be tested with tools such as:

\begin{lstlisting}
AddressSanitizer
UndefinedBehaviorSanitizer
Valgrind
\end{lstlisting}

Example development flags:

\begin{lstlisting}[language=make]
CFLAGS += -fsanitize=address,undefined
LDFLAGS += -fsanitize=address,undefined
\end{lstlisting}

These need not be production flags, but allocation, buffer, parser, and
ownership code should be tested with such tools when practical.

\section{27. Comments and public API
documentation}\label{27-comments-and-public-api-documentation}

Comments should explain why, ownership, protocol requirements, hardware
constraints, or non-obvious lifetime assumptions.

For public functions, ownership documentation is part of the API.

Comments shoule be done with // and not /* ... */.


Example:

\begin{lstlisting}[language=C]
//----------------------------------------------------------------------
// rx_device_name_copy:
//
// Returns a newly allocated copy of the device name.
//
//  Ownership:
//  The caller owns the returned string and shall release it with free().
//
// Returns:
//   Newly allocated string on success, NULL on failure.
//----------------------------------------------------------------------
char *
rx_device_name_copy(const rx_device_t *device);
\end{lstlisting}

\section{28. Example}\label{28-example}

\begin{lstlisting}[language=C]
#include "rx-message.h"

#include <stdlib.h>
#include <string.h>

rx_message_t *
rx_message_new(
    const void *data,
    size_t data_size
)
{
    rx_message_t *message;

    message = calloc(1, sizeof(*message));

    if (message == NULL)
    {
        return NULL;
    }

    if (data_size > 0)
    {
        message->data = malloc(data_size);

        if (message->data == NULL)
        {
            free(message);
            return NULL;
        }

        memcpy(
            message->data,
            data,
            data_size
        );

        message->data_size = data_size;
    }

    return message;
}

void
rx_message_free(rx_message_t *message)
{
    if (message == NULL)
    {
        return;
    }

    free(message->data);
    free(message);
}
\end{lstlisting}

Singel line comment is like

\begin{lstlisting}[language=C]
free(message->data); // Free memory
\end{lstlisting}


And for start of lager block there should be 70 hypens like

\begin{lstlisting}[language=C]
//----------------------------------------------------------------------
//  Comment
//----------------------------------------------------------------------
\end{lstlisting}
  


\section{29. Review checklist}\label{29-review-checklist}

Before committing C code, check:

\begin{enumerate}
\def\labelenumi{\arabic{enumi}.}
\tightlist
\item
  Files and modules are named consistently.
\item
  Identifiers are lowercase with underscores.
\item
  Public symbols carry an appropriate prefix.
\item
  Names are descriptive and single-letter abbreviations are avoided.
\item
  Braces and four-space indentation are consistent.
\item
  Every pointer lifetime is understandable.
\item
  Every allocation has a matching release.
\item
  All failure paths are leak-free.
\item
  \passthrough{\lstinline!realloc()!} is used safely.
\item
  Buffer lengths are explicit.
\item
  Unsafe string functions are avoided.
\item
  Protocol lengths and integer widths are validated.
\item
  Return values are checked.
\item
  All acquired resources are released in every path.
\item
  Public ownership rules are documented.
\item
  The normal build has no unexplained warnings.
\item
  Memory-sensitive code has been checked with sanitizers or Valgrind
  when practical.
\item
  The implementation is simpler than a more abstract alternative.
\end{enumerate}

\section{30. Core rule}\label{30-core-rule}

\begin{quote}
Ownership, lifetime, bounds, and cleanup shall be visible in the code.
\end{quote}

C code is not complete merely because the successful path works. The
cleanup and failure paths are part of the implementation.


\end{document}
